\section{Síntesis y ampliación}

\subsection{Conclusiones centrales}

\begin{enumerate}
\item Miaoya (妙鸭) es un producto de internet que aprovecha muy bien las capacidades de la IA, pero no es AI Native en el sentido que Zhang Yueguang (张月光) definió después.
\item El cambio decisivo de AI Native es pasar del diseño de flujos al diseño del contexto: el gerente de producto debe organizar la información, las herramientas y la retroalimentación que el modelo puede ver.
\item One Way Door es una señal importante para reconocer un producto acertado a largo plazo: una vez que el usuario migra, ya no quiere volver a la solución anterior.
\item La organización de un producto de IA deja de ser un trabajo lineal y pasa a tener dos etapas: primero se validan las capacidades del modelo y el context; después se encapsulan el flujo y la experiencia.
\item Zhang Yueguang entiende esta generación de IA como la creación de una «población de IA», y no solo como la creación de servicios de IA.
\item La división Will / Skill indica que la persona se encarga de los objetivos, el juicio y el gusto, y la IA, de la ejecución y de la ampliación de capacidades.
\item Xingmian (星眠) explora el contenido y la compañía con IA; Dokie explora un Agent de creación y expresión, y la superación de los límites de las capacidades del usuario.
\item Un producto de Agent puede empezar por una ruta de retroalimentación corta y frecuente, en lugar de perseguir solamente tareas fuera de línea de muy largo alcance.
\item El LM sigue siendo el núcleo de los Agent actuales, pero esa es una cuestión distinta de si la inteligencia en el mundo físico necesita ir más allá del lenguaje.
\end{enumerate}

\subsection{Preguntas abiertas}

Se dejan preguntas abiertas al final porque este episodio ofrece sobre todo metodología de producto y juicios de emprendimiento, no conclusiones que el mercado ya haya validado por completo. Las siguientes preguntas determinarán si Dokie y los productos de Agent similares pueden pasar de ser herramientas a ser puntos de entrada.

\begin{itemize}
\item ¿Los productos AI Native formarán una metodología estable, o cada tipo de producto tendrá que reinventar el context design?
\item ¿El mejor punto de entrada para un Agent es la creación y expresión, el procesamiento de información, el código, la automatización de oficina o un amigo de IA con personalidad?
\item ¿Pueden los juegos otome con IA convertirse en una categoría grande, o solo serán un camino estrecho dentro del contenido con IA?
\item ¿Cuánto tiempo podrá el LM seguir siendo el núcleo de los Agent? ¿Cambiarán este juicio los modelos multimodales?
\item ¿Los usuarios pagarán por un «amigo de IA», o solo por una «herramienta que trabaje de forma confiable»?
\end{itemize}

\subsection{Lecturas adicionales}

\begin{itemize}
\item EP139, historia técnica de los Agent: para entender la trayectoria de los Agent desde la invocación de herramientas hasta su alcance social.
\item EP135, Tristan / Elys: para entender los dobles digitales, el flujo del context y las redes sociales de IA.
\item EP113, entrevista a Yang Zhilin (杨植麟): Agentic LLM, el cerebro en una cubeta y «El comienzo del infinito».
\item Casos de coding agent como Cursor y Claude Code: para entender los productos One Way Door.
\item Materiales sobre juegos con IA, narrativa interactiva, Live2D y compañía virtual: para entender la línea de contenido de Xingmian.
\end{itemize}

\begin{importantbox}{El juicio final}
Lo que más vale llevarse de este episodio no es el juicio contraintuitivo de que «Miaoya no es AI Native», sino la metodología que lo sustenta: cuando tanto la entrada como la salida son abiertas, el producto deja de ser solo un flujo; se convierte en un sistema formado en conjunto por el contexto, el modelo, las herramientas, la retroalimentación, el gusto y la percepción del usuario.
\end{importantbox}

\end{document}
